%% kapa: exam=1 round=24 year=2013 period=2 form=A subject=accounting subject_ko=회계학 seq=3 number=3
%% source: Q-Net 감정평가사 기출문제 (cst00309) artlSeq=5073802 file=extracted/1008/5073802/2013년 제24회 감정평가사 1차시험 기출문제/2교시(A형).hwp
%% conversion: kordoc 4.24.51 -> tools/gen1.py (auto); review=proofread-sonnet; flags=none
\begin{question}{3}
(주)한국은 20×1년 1월 1일 내용연수 10년, 잔존가치 ￦0인 건물을 ￦20,000,000에 취득하여 정액법으로 상각하고 있다. 원가모형을 적용하는 건물의 20×2년 12월 31일 회수가능액은 ￦12,800,000이고, (주)한국은 이에 대한 손상차손을 인식하였다. 20×3년 12월 31일에는 회수가능액이 ￦16,000,000으로 회복되는 경우, 건물과 관련된 회계처리가 (주)한국의 20×3년 당기순이익에 미치는 영향은?
\begin{choices}
  \item ￦1,200,000 증가
  \item ￦1,200,000 감소
  \item ￦3,200,000 증가
  \item ￦3,200,000 감소
  \item ￦4,200,000 증가
\end{choices}
\end{question}
\answer{1}
